% =====================================================================
% colors.tex　语义化颜色层（全书颜色的唯一来源）
%
% 规则
%   1. 正文、栏目框和图中只使用本文件定义的颜色名，以及中性色 black、white、gray!N；
%      不要在别处直接写 red!4、orange!65!black 这样的具体颜色.
%   2. 需要浅色或深色变体时，在语义色上调配，例如 accent!8、secondary!6、series1!25.
%   3. 改配色只改本文件；印刷黑白版（\notemode = printbw）的灰度映射也只在这里处理.
%
% 依赖：xcolor（elegantbook 已加载）、tcolorbox（document.tex 导言区加载）.
% 由 document.tex 在“全书统一记号”之前 \input.
% =====================================================================

% ---------- 1. 基础色 ----------
% 主色：封面色带、标题页色带、课程结构图、题词框
\definecolor{customcolor}{RGB}{104,200,240}   % 旧名，模板与封面仍在用，勿删
\colorlet{accent}{customcolor}
\colorlet{accentdark}{accent!75!black}
\colorlet{coverlinecolor}{accent}             % 模板：标题页的色带

% 辅助色：只作为调配的“原料”，正文中优先用下面的具体角色名
\colorlet{secondary}{orange}   % 第二部分底色、专业连线、映射的到达集合
\colorlet{tertiary}{violet}    % 第三部分底色、建模实验
\colorlet{alert}{red}          % 常见误区
\colorlet{success}{green}      % 计算实践

% ---------- 2. 图中的类别色 ----------
% 用来区分图中不同的对象（向量、线路、曲线……），本身没有先后或褒贬含义.
% 黑白印刷时颜色会失效，因此图中必须同时用线型、形状或文字区分（见 ROADMAP v0.14）.
\colorlet{series1}{accentdark}
\colorlet{series2}{secondary!85!black}
\colorlet{series3}{alert!70!black}
\colorlet{series4}{success!45!black}
\colorlet{series5}{tertiary!70!black}
% 旧名：正文中已大量使用，保留为别名，换用线型样式时再逐步替换
\colorlet{figblue}{series1}
\colorlet{figorange}{series2}
\colorlet{figred}{series3}
\colorlet{figgreen}{series4}
\colorlet{figviolet}{series5}

% ---------- 3. 贯穿全书的四个建模案例（课程结构图中的标记） ----------
\colorlet{caseA}{series1}   % 投入产出
\colorlet{caseB}{series2}   % 物流网络
\colorlet{caseC}{series3}   % 判例引证与法条检索
\colorlet{caseD}{series4}   % 金融

% ---------- 4. 栏目框 ----------
% 每类栏目一对颜色：<名>back 底色，<名>frame 边框与标题栏
\colorlet{motiveback}{accent!6}          \colorlet{motiveframe}{accent!70!black}      % 思维动机
\colorlet{linkback}{secondary!5}         \colorlet{linkframe}{secondary!65!black}     % 专业连线
\colorlet{historyback}{gray!5}           \colorlet{historyframe}{gray!60!black}       % 历史回响
\colorlet{pitfallback}{alert!4}          \colorlet{pitfallframe}{alert!55!black}      % 常见误区
\colorlet{codeback}{success!3}           \colorlet{codeframe}{success!45!black}       % 计算实践
\colorlet{modelback}{tertiary!4}         \colorlet{modelframe}{tertiary!60!black}     % 建模实验
% 章首题词框（sidefig.tex）
\definecolor{quotebg}{RGB}{235,243,251}
\colorlet{epigraphframe}{accent}

% 六类栏目共用的基础样式：ln box={<颜色前缀>}{<标题>}
\tcbset{
	ln box/.style 2 args={
		breakable, enhanced,
		colback=#1back, colframe=#1frame,
		fonttitle=\bfseries\sffamily, title={#2},
	},
}

% ---------- 5. 输出模式 ----------
% printbw：全部颜色按亮度转为灰度（xcolor 的目标颜色模型）.
% 这里只搭框架：灰阶层次、栏目区分、图片处理留到 v0.14、v0.15 细化.
\ifbwversion
	\selectcolormodel{gray}
\fi
